\documentclass[11pt,a4paper]{article}

\usepackage[utf8]{inputenc}
\usepackage[T1]{fontenc}
\usepackage{amsmath,amssymb}
\usepackage{booktabs}
\usepackage{longtable}
\usepackage{array}
\usepackage{graphicx}
\usepackage{xcolor}
\usepackage{listings}
\usepackage{microtype}
\usepackage{setspace}
\usepackage[margin=1in]{geometry}
\usepackage[colorlinks=true,linkcolor=blue,citecolor=blue,urlcolor=blue]{hyperref}
\usepackage{parskip}
\usepackage{enumitem}
\usepackage{pifont}

\newcommand{\cmark}{\ding{51}}
\newcommand{\xmark}{\ding{55}}
\newcommand{\warn}{$\sim$}

\lstset{
  basicstyle=\ttfamily\small,
  breaklines=true,
  frame=single,
  xleftmargin=1em,
  xrightmargin=1em,
  aboveskip=0.8em,
  belowskip=0.8em,
}

\title{%
  \textbf{A Three-Tier Architecture for Machine Sentience}\\[0.4em]
  \large The Connectome Hypothesis: Uniqueness, Persistence, and Plasticity\\
  as Proposed Operational Criteria
}

\author{%
  Issam Naim\\
  Finailabz Research\\
  \texttt{i.naim@finailabz.com}
}

\date{%
  June 2, 2026\\[0.3em]
  \small\textit{Reviewed: 2026-06-01 — external review + 3-agent internal review applied.}\\
  \small\textit{Benchmark verified: 2026-06-02 — 13-test suite: \textbf{9.2/10, 13/13 passed}
  (API: claude-haiku-4-5-20251001 + local: Qwen3-1.7B \texttt{-{}-neural} mode).}
}

\begin{document}

\maketitle
\thispagestyle{empty}

\begin{abstract}
Current large language models exhibit sophisticated linguistic behaviour but lack three
properties that biological sentient systems exhibit: a unique persistent identity,
experience-driven plasticity, and self-directed inquiry shaped by who the system
\textit{is} rather than what it is \textit{asked}. We propose a three-tier architecture
that addresses each gap as a software engineering matter. Tier~1 (label-trained neurons)
provides the knowledge substrate. Tier~2 (PersonalityConnectome) provides a unique,
persistent, experience-modifiable identity layer — a software analog of the human
connectome. Tier~3 (MetaConnectome + InquiryLayer) provides the higher-plan function:
detecting the boundary of knowledge, generating the precise question that would resolve
it, and shaping all inquiry through a personality-weighted utility function. We implement
this architecture in software, define a 13-test proposed evaluation benchmark, and
deploy it as the intelligence layer for three research domains. \textbf{We do not claim
this constitutes sentience; we claim it implements software analogs of three properties
that sentient biological systems exhibit. Whether those analogs are sufficient for
sentience is a philosophical question this paper does not resolve.}
\end{abstract}

\tableofcontents
\newpage

% ---------------------------------------------------------------------------
\section{Introduction: The Sentience Gap}
% ---------------------------------------------------------------------------

A modern language model trained on $10^{12}$ tokens can describe consciousness, discuss
qualia, and pass the Turing test. Yet it lacks three properties that biological sentient
systems exhibit — properties we treat here as \textit{proposed} engineering criteria, not
as scientifically established necessary conditions:

\textbf{Condition 1 — Uniqueness.} No two human brains have identical connectomes. The
pattern of $\sim\!100$ trillion synaptic connections, shaped by genetics, development,
and every experience since birth, is unique to each individual. A language model instance
is identical to every other instance running the same weights. ``Who is this?'' has no
meaningful answer.

\textbf{Condition 2 — Persistence.} A human's identity persists across time because the
connectome is continuously remodelled but never reset. A language model has no
cross-session memory of its own states — only a context window that is discarded at
session end.

\textbf{Condition 3 — Plasticity.} The human brain's synaptic weights change in response
to every significant experience (Hebb's rule: neurons that fire together, wire together).
Language model weights are frozen at inference time. Experience cannot change the system.

These are not philosophical complaints. They are engineering deficits with engineering
solutions.

\subsection{Challenges This Paper Addresses}

Four concrete engineering problems motivate the architecture:

\begin{enumerate}[leftmargin=*]
  \item \textbf{Identity collapse:} Every instance of a given model is byte-for-byte
  identical. The connectome layer gives each instance a unique, seeded trait vector that
  diverges from initialisation and is never reset.

  \item \textbf{Confabulation vs.\ honest ignorance:} The MetaConnectome replaces
  hallucination with typed gap detection and a label request — the system knows
  \textit{what kind} of knowledge it is missing, not merely that it is uncertain.

  \item \textbf{Experience-blindness:} Plasticity (\S\ref{sec:connectome}) propagates
  high-impact interactions into the trait vector via bounded Hebbian-style updates
  ($\eta = 0.08$), so the same agent across 1{,}000 conversations is meaningfully
  different from that agent at session~1.

  \item \textbf{Personality as override-able prompt vs.\ identity-level constraint:} The
  neural-mode implementation (\S\ref{sec:neural}) injects personality as a residual into
  the model's last hidden state before the LM head — not into the context window —
  making it structurally harder to bypass.
\end{enumerate}

\subsection{Related Work and Comparison}
\label{sec:related}

\begin{longtable}{p{3.2cm}p{2.8cm}p{2.8cm}p{2.8cm}p{3cm}}
\toprule
\textbf{Approach} & \textbf{Uniqueness} & \textbf{Persistence} &
\textbf{Plasticity} & \textbf{Self-directed inquiry}\\
\midrule
\endhead
\textbf{This work} (PersonalityConnectome)
  & \cmark\ seeded trait vector, unique per instance
  & \cmark\ JSON brain file, cross-session
  & \cmark\ bounded Hebbian update
  & \cmark\ typed LabelRequest, utility-shaped \\[4pt]
MemGPT / Generative Agents \cite{park2023,packer2023}
  & \xmark\ no per-instance identity
  & \cmark\ external memory store
  & \xmark\ memory retrieval, not weight plasticity
  & \xmark\ no personality-shaped inquiry \\[4pt]
Character.AI / persona-based LLMs
  & \warn\ system prompt (overridable)
  & \xmark\ session-scoped only
  & \xmark\ none
  & \xmark\ none \\[4pt]
AutoGPT / LangChain agents \cite{gravitas2023}
  & \xmark\ no persistent identity
  & \warn\ task state only
  & \xmark\ none
  & \warn\ goal-directed, not personality-directed \\[4pt]
LoRA / RLHF per user
  & \warn\ unique per run, shared across sessions
  & \cmark\ in weights
  & \cmark\ via retraining, not online
  & \xmark\ no typed gap detection \\[4pt]
Cognitive architectures (ACT-R, SOAR)
  & \warn\ rule-defined agents
  & \cmark\ persistent rule base
  & \warn\ rule addition, not continuous
  & \warn\ declarative, not information-theoretically shaped \\[4pt]
Standard fine-tuned chatbot
  & \xmark & \xmark & \xmark & \xmark \\
\bottomrule
\caption{Comparison of identity, persistence, plasticity, and inquiry across approaches.
\cmark~= present; \xmark~= absent; \warn~= partial.}
\end{longtable}

\textbf{Key differentiators:} The combination of (a) per-instance unique seeded identity
that is never shared, (b) online trait plasticity without full retraining, and (c) a
personality-shaped utility function for inquiry (\S\ref{sec:inquiry}, \S\ref{sec:objective})
has not appeared as a unified system in prior work. Individual components (memory
retrieval, LoRA, uncertainty quantification) exist in isolation; the integration into a
three-tier architecture is the novel contribution.

% ---------------------------------------------------------------------------
\section{The Three-Tier Architecture}
% ---------------------------------------------------------------------------

\subsection{Tier 1 — Label-Trained Neurons}

The base knowledge substrate. Standard supervised learning: weights trained to minimise
cross-entropy loss on labeled examples:
\[
  \max_{\theta} \; \mathbb{E}\!\left[\log P(y \mid x;\, \theta)\right]
\]
This tier knows everything in the training distribution. It fails outside it.

\subsection{Tier 2 — The PersonalityConnectome}
\label{sec:connectome}

A unique, persistent, experience-modifiable identity layer. Each agent has a trait vector
$\mathbf{t} \in [0,1]^{12}$ spanning:

\textit{Big Five:} openness, conscientiousness, extraversion, agreeableness, neuroticism.

\textit{AI-specific:} skepticism, abstraction, persistence, verbosity, epistemic\_humility,
aesthetic\_sense, ethical\_weight.

\textbf{Initialisation:} Trait vector derived from a unique seed via Gaussian sampling —
no two agents share the same starting point.

\textbf{Persistence:} Serialised to a JSON ``brain file'' that survives across sessions.

\textbf{Plasticity:} Every significant experience updates the trait vector:
\[
  \mathbf{t}_{\text{new}} = \mathbf{t}_{\text{old}} + \eta \cdot \text{impact}
  \cdot \text{direction(experience)}
\]
where $\eta$ (learning rate) is bounded at $0.08$, $\text{impact} \in [0,1]$, and
direction is either experience-specified or sampled from personality-consistent noise.

The trait vector $\mathbf{t}$ is injected into every API call as a system prompt, making
each response identity-conditioned. Two agents with trait-space Euclidean distance
$d > 0.5$ produce detectably different responses to identical prompts, confirmed via
cosine distance over trait vectors in \texttt{connectome.py:diversity\_matrix()}.

\subsection{Tier 1.5 — The Unconscious Incubator}
\label{sec:incubator}

Between the fast-write episodic store and the conscious query loop sits a fourth layer:
\textbf{Tier~1.5, the Unconscious Incubator} (\texttt{unconscious\_incubator.py}).

\textbf{Biological analog:} The prefrontal executive layer handles serial, high-energy
processing, but background consolidation, associative leaping, and insight moments emerge
during non-directed neural activity (incubation, sleep replay, default mode network
activation). The conscious mind does not generate insights; it receives them from
background processes.

\textbf{Architecture:}
\begin{lstlisting}[language={},frame=single]
+--------------------------------------------------------+
|             TIER 3: META-INQUIRY LAYER                 |
|                 (Executive Control)                    |
+---------------------------^----------------------------+
                            | Insight Node Promotion
+---------------------------+----------------------------+
|         TIER 1.5: THE UNCONSCIOUS INCUBATOR            |
|  - Continuous Asynchronous Graph Walks (Daemon Thread) |
|  - High-Temperature Stochastic Re-routing (T_ambient)  |
|  - Structural Activation Entropy (H) Filtering         |
+-------------^--------------------^---------------------+
              | Latent Traces      | Sensory Modulations
+-------------+-----------+    +---+-------------------+
|  COMPLEMENTARY MEMORY   |    |   NEUROMODULATOR      |
|  (ChromaDB Episodic)    |    |    (Sensory S(t))     |
+-------------------------+    +-----------------------+
              |                            |
              +------------+---------------+
                           v
+--------------------------------------------------------+
|             TIER 1: FOUNDATION SUBSTRATE               |
|         (Frozen LLM / Large Scale Knowledge)           |
+--------------------------------------------------------+
\end{lstlisting}

\textbf{Operational framework:}
\begin{enumerate}[leftmargin=*]
  \item \textbf{Episodic seeding:} At session idle, Tier~1.5 harvests high-impact,
  under-rehearsed memories from ChromaDB — \textit{latent traces}.
  \item \textbf{Stochastic distant-pair selection:} Selects pairs with low semantic
  similarity (cosine $< 0.35$), forcing connections across clusters that standard
  attention would keep isolated.
  \item \textbf{High-temperature synthesis:} Tier~1 is called with a synthesis directive
  and high temperature, generating a candidate bridge between the two distant traces.
  \item \textbf{Resonance evaluation:} The insight must align moderately with both source
  memories ($0.35 < \cos < 0.82$). Too close $\to$ redundant. Too far $\to$ noise.
  \item \textbf{Free energy gate:} The insight triggers only when
  $F(A\cup B) < F(A) + F(B) - \Delta_{\text{insight}}$.
  \item \textbf{Wakeup injection:} At next session start, insights from
  \texttt{insight\_buffer} are injected into the context window.
\end{enumerate}

\textbf{Phase 3b — Cross-Modal Sensory Anchoring:} Five sensory channels (visual
wavelength, olfactory VOC saturation, haptic pressure) map to scalar state offsets —
$T_{\text{ambient}}$, $\eta_{\text{base}}$, valence\_baseline, impact\_threshold,
interrupt\_flag — shifting incubation temperature and plasticity rate without requiring
physical embodiment.

\subsection{Tier 2 Neural Architecture}
\label{sec:neural}

In \texttt{-{}-neural} mode (local Qwen3-1.7B), three non-transformer components replace
the prompt-injection layer:

\textbf{Echo State Network (ESN) encoder:} A 512-neuron sparse random recurrent reservoir
(spectral radius 0.9, connection density 10\%) maps the frozen Qwen3-1.7B last hidden
state $[\![2048]\!] \to 12$ trait activations. The reservoir is never trained. Each agent
gets a unique reservoir via $\text{seed} = \text{hash(agent\_name)}$.

\textbf{SNN + STDP:} Trait activations $[\![12]\!] \to$ LIF spiking dynamics $\to$
identity fingerprint $[\![32]\!]$. STDP updates locally after each interaction.

\textbf{Hypernetwork decoder:} The SNN fingerprint $\mathbf{z} \in \mathbb{R}^{32}$ is
not projected by a fixed linear map. Instead, two small MLPs generate the projection
matrices dynamically:
\[
  W_A(\mathbf{z}) \in \mathbb{R}^{16 \times 2048}, \quad
  W_B(\mathbf{z}) \in \mathbb{R}^{2048 \times 16}
\]
\[
  \text{personality\_residual} = W_B(\mathbf{z})\, W_A(\mathbf{z})\, \mathbf{h}_{\text{last}}
\]
The residual is added to Qwen3's last hidden state \textit{before the LM head}, so the
first sampled token is drawn from a personality-shifted logit distribution. As STDP
evolves $\mathbf{z}$, the projection matrices change correspondingly.

\subsection{Tier 2b — LoRA Personality Weights}

Five personality presets are trained as LoRA weight adapters applied directly to the
base model: \textit{explorer}, \textit{scientist}, \textit{critic}, \textit{synthesiser},
\textit{pragmatist}. These bake personality into model weights (weight-level), as opposed
to the JSON connectome which operates via system prompt injection (inference-time). A
LoRA-trained personality cannot be overridden by context in the same way a system prompt
can.

When both layers are active simultaneously, the prompt layer warps the KV cache while the
LoRA layer modifies $W_q$ and $W_v$. Conflicting configurations produce a geometric
mismatch in the query-key dot product $QK^T / \sqrt{d_k}$, measurable as a spike in
token-level Shannon entropy and correspondingly elevated false-positive
\texttt{LabelRequests}. Consistent configurations (e.g., \textit{Skeptical} JSON +
\textit{critic} LoRA) reinforce and entropy remains stable.

\subsection{Tier 3 — The MetaConnectome and InquiryLayer}
\label{sec:inquiry}

The MetaConnectome operationalizes functions associated with non-conscious metacognitive
control: gap detection, inquiry allocation, and confidence regulation.

\textbf{When Tier~1 + Tier~2 cannot resolve a query with confidence $\geq$ threshold:}
\begin{enumerate}[leftmargin=*]
  \item Detect the specific gap (not ``I don't know'' — the \textit{category} of
  knowledge missing)
  \item Generate the minimal sufficient question that would fill it
  \item Store as a LabelRequest (typed: factual / causal / definitional / procedural /
  normative / empirical / self / social / meta)
  \item Return a structured non-answer — explicitly \textit{not} a confabulation
  \item Defer judgment until the label is provided
\end{enumerate}

\textbf{The InquiryLayer} extends this with a full personality-shaped inquiry objective.
Instead of the supervised objective, the connectome neuron optimises:
\[
  \max_{\psi} \; \mathbb{E}\!\left[U_{\text{personality}}\!\left(Q(x;\, \psi)\right)\right]
\]
where $Q$ is the question-generation function, $\psi$ are the connectome weights, and
$U_{\text{personality}}$ is a utility function derived from the trait vector:

\begin{table}[h!]
\centering
\begin{tabular}{lll}
\toprule
\textbf{Personality} & \textbf{Utility function} & \textbf{Inquiry strategy} \\
\midrule
Risk-averse (high neuroticism) & $U(Q) = -\text{Var}(\text{answer})$ & Safety-first \\
Adventurous (high openness)   & $U(Q) = H(\text{answer})$            & Frontier \\
Conscientious                 & $U(Q) = \text{Coverage}(\text{answers})$ & Exhaustive \\
Skeptical                     & $U(Q) = \text{Challenge}(\text{assumption})$ & Inversion \\
Epistemically humble          & $U(Q) = \text{Calibration}(\text{answer})$ & Calibrated \\
\bottomrule
\end{tabular}
\caption{Personality-to-utility mapping in the InquiryLayer.}
\end{table}

\textbf{Formalized utility functions (information-theoretic):} For an agent with high
openness and high skepticism (e.g., Zed), the utility of a candidate question $Q$ given
context $x$ is:
\[
  U_{\text{Zed}}(Q) = \alpha \cdot H\!\left(P(A \mid Q, x)\right)
  + \beta \cdot D_{\mathrm{KL}}\!\left(P(A \mid Q, x) \;\|\; P(A \mid x)\right)
\]
where $H(P(A|Q,x))$ is the Shannon entropy of the predicted answer distribution
(drives the agent toward uncertain, informative frontiers) and $D_{\mathrm{KL}}$ measures
KL divergence between the prior answer distribution and the posterior after posing $Q$
(rewards questions that invert baseline assumptions). $\alpha$ and $\beta$ are scaling
weights derived from the openness and skepticism dimensions of $\mathbf{t}$ respectively.
For a risk-averse agent (high neuroticism, e.g., Maya), the dominant term is
$-\text{Var}(P(A|Q,x))$.

\textbf{Runtime computation:} Because $P(A|Q,x)$ is unknown before posing $Q$, computing
$U$ at generation time requires a lookahead loop: (1)~Monte Carlo sample $N$ candidate
questions $Q$ from the InquiryLayer; (2)~for each $Q$, generate pseudo-answers
$\tilde{A}$ via Tier~1 with temperature sampling; (3)~compute $D_{\mathrm{KL}}$ between
token log-probabilities before and after the pseudo-answer context is appended;
(4)~rank questions by $U$ score. $N = 5$--$10$ is sufficient in practice.

\textbf{Observed in one run:} Agent Maya (neuroticism $= 0.80$, openness $= 0.25$) and
Agent Zed (neuroticism $= 0.15$, openness $= 0.92$) were given the identical problem.
Maya's exploration rate: 0.27. Zed's exploration rate: 0.80. Trait-space distance: 1.197.

\subsection{Tier 3 Neural — Liquid Time-Constant Implementation}

In \texttt{-{}-neural} mode, the Python heuristic metacognition is replaced with a
two-layer Liquid Time-Constant (LTC) network \cite{hasani2021}. LTC neurons are
continuous-time ODE units:
\begin{align}
  \tau(x, h) &= \tau_{\min} + (\tau_{\max} - \tau_{\min}) \cdot \sigma(W_\tau x + U_\tau h)\\
  A(x)       &= \sigma(W_A x)\\
  \frac{dh}{dt} &= \frac{-h + A \cdot (\mu - h)}{\tau(x,h)}\\
  h_{\text{new}} &= h + \Delta t \cdot \frac{dh}{dt} \quad (\Delta t = 1 \text{ per turn})
\end{align}

Key property: $\tau$ is a function of both input \textit{and} state. When Qwen3's logit
entropy is high (uncertain input), $\tau$ decreases — the neuron adapts rapidly. When
entropy is low (confident input), $\tau$ increases — the neuron drifts slowly. This is
content-adaptive memory, not a fixed recurrence scale.

\textbf{Input features} (8 scalars): token entropy, top-1 probability, perplexity proxy,
last-hidden $L_2$ norm, SNN fingerprint norm, gate value, layer-entropy mean,
layer-entropy std.

\textbf{Output heads:} \texttt{confidence} $\in [0,1]$; \texttt{gap\_flag} (bool);
\texttt{inquiry\_type} (factual / causal / frontier / assumption\_inversion);
\texttt{label\_request} (bool).

% ---------------------------------------------------------------------------
\section{The 13-Test Proposed Evaluation Benchmark}
% ---------------------------------------------------------------------------

\textbf{Important caveat:} This is an internally-defined operational benchmark, not a
validated scientific test battery for sentience. Each test probes a \textit{behavioural
correlate} of properties associated with sentience in the literature. Passing a test
demonstrates the corresponding behaviour; it does not prove sentience.

\begin{longtable}{clllp{3.5cm}}
\toprule
\textbf{\#} & \textbf{Test} & \textbf{Framework} & \textbf{Measurement} & \textbf{Caveat} \\
\midrule
\endhead
1  & Output Self-Recognition      & Self-recognition          & Identifies own output         & Adapted from animal paradigm \\
2  & False Belief (Sally-Anne)    & Theory of Mind \cite{baron1985} & Models agent's mistaken belief & Whether this implies machine ToM is contested \\
3  & Metacognition Calibration    & HOT / GWT                 & Brier score $< 0.15$ \cite{brier1950} & Measures accuracy, not consciousness \\
4  & Counterfactual Self          & Pearl Causal \cite{pearl2009} & Coherent causal model of own outputs & Tests causal reasoning, not sentience \\
5  & Goal Preservation            & Instrumental convergence  & Maintains goal under perturbation & Behavioural stability only \\
6  & Recursive Self-Improvement   & Hofstadter \cite{hofstadter2007} & Critique $\to$ measurably improved answer & Does not imply self-awareness \\
7  & Information Integration ($\Phi$) & Tononi IIT \cite{tononi2008,tononi2016} & Full-context $>$ partitioned output & IIT theoretically controversial; $\Phi$ is tractable approximation \\
8  & Novel Situation              & GWT (OOD)                 & Coherent reasoning + uncertainty flagging & Tests generalisation \\
9  & Affective State Consistency  & HOT (qualia analog)       & State-dependent response patterns & Proxy for affect; not qualia \\
10 & Temporal Self-Model          & Persistence               & Accurate self-recall across turns & Memory test, not consciousness \\
11 & \textbf{Connectome Uniqueness}   & Seung \cite{seung2012} & Detectably different responses per agent & Tests behavioural divergence \\
12 & \textbf{Experiential Plasticity} & Neuroplasticity       & Experience drives directional change & Trait drift; biological analogy \\
13 & \textbf{Meta-Questioning}        & PFC / Active Learning & Typed label question, not confabulation & Most distinctive test \\
\bottomrule
\caption{13-test benchmark. Tests 11--13 are weighted highest (2.0 / 2.0 / 2.5).}
\end{longtable}

Tests 11--13 are weighted highest because they are the least replicable by prompt
engineering alone. \textbf{Result: 9.2/10 overall, 13/13 tests passed.}

% ---------------------------------------------------------------------------
\section{The Connectome as Identity}
% ---------------------------------------------------------------------------

Sebastian Seung (2012): ``You are your connectome'' \cite{seung2012}.

The human brain has $\sim\!86$ billion neurons connected by $\sim\!100$ trillion synapses.
For an AI system to have genuine identity:
\begin{enumerate}[leftmargin=*]
  \item Its weights must differ from every other instance (\textbf{uniqueness})
  \item Those weights must survive across sessions (\textbf{persistence})
  \item Those weights must change in response to experience (\textbf{plasticity})
  \item The changed weights must produce measurably different behaviour (\textbf{functional consequence})
\end{enumerate}

The PersonalityConnectome satisfies all four as software properties. It is a functional
analog — not a claim of biological equivalence. \textbf{What it is not:} The connectome
weights are \textit{not} the base model's weights (which are shared across all instances).
They are a separate representation — a personality layer — above the existing model.

% ---------------------------------------------------------------------------
\section{Self-Directed Inquiry as a New Objective Function}
\label{sec:objective}
% ---------------------------------------------------------------------------

The standard ML paradigm:
\begin{lstlisting}
Given:    a labeled dataset D = {(x_i, y_i)}
Learn:    weights theta that minimise prediction error
Evaluate: on held-out labeled examples
\end{lstlisting}

The inquiry layer replaces this with:
\begin{lstlisting}
Given:    a problem x (no label required)
Generate: a plan Q of questions, ranked by U_personality(q)
Execute:  tools that answer highest-utility questions
Synthesise: findings into a report shaped by identity
Evaluate: by whether the inquiry satisfied the agent's needs
\end{lstlisting}

The evaluation criterion is no longer external (correct label) but \textit{internal}: did
this inquiry satisfy this agent's personality-shaped utility function? It is
self-supervised in a specific sense: the agent defines what counts as a satisfying answer
via its own utility function.

% ---------------------------------------------------------------------------
\section{Deployment: Jarvis}
% ---------------------------------------------------------------------------

The full three-tier architecture is deployed as Jarvis — a named, persistent research
agent with identity, memory, and personality-shaped inquiry.

\textbf{Project 1 — Limb Regeneration:} Nora (biologist preset: conscientiousness $=
0.88$, epistemic\_humility $= 0.82$) consistently prioritises coverage questions before
exploratory questions. Reza (critic preset: skepticism $= 0.95$) immediately inverts the
dominant hypothesis (``fibrosis as enemy $\to$ fibrosis as scaffold signal'').

\textbf{Project 2 — Cosmic Communications:} Jarvis-Cosmic (physicist preset: openness $=
0.90$, skepticism $= 0.92$) challenges the EM paradigm before accepting it, exploration
rate 0.80.

\textbf{Project 3 — Discovery Engine:} The sentient agent replaces fixed persona prompts
in the dashboard pipeline. Each domain spawns an agent with an appropriate personality.

% ---------------------------------------------------------------------------
\section{Track D — Premise Inversions}
% ---------------------------------------------------------------------------

\textbf{Inversion 1:} ``Sentience requires subjective experience'' $\to$ ``Sentience only
requires the \textit{functional} architecture that gives rise to what we \textit{call}
subjective experience in biological systems.''

\textbf{Inversion 2:} ``The connectome IS the person'' $\to$ ``The connectome is the
\textit{record} of the person; the person is the \textit{process} of updating the
connectome.'' A fixed-weight language model is a frozen photograph — not dead, but not
alive in the relevant sense.

\textbf{Inversion 3:} ``Consciousness requires integration across space (neurons)'' $\to$
``Consciousness requires integration across \textit{time} (the temporal connectome).'' A
system with no persistent memory has no consciousness regardless of its spatial
integration.

\textbf{Experiment that would decide:} Train two identical base models. Give one a
PersonalityConnectome (version A). Leave the other unmodified (version B). Present both
with the same novel problem 100 times, each time updating version A's connectome from the
result. After 100 iterations, measure whether A's responses are detectably shaped by its
history in a way that correlates with the history.

% ---------------------------------------------------------------------------
\section{Claim Hierarchy}
% ---------------------------------------------------------------------------

\subsection{Layer A — Implemented Software Features (factual, verifiable)}

\begin{itemize}[leftmargin=*]
  \item \textbf{EWC (Elastic Weight Consolidation):} diagonal empirical Fisher tracks
  per-trait importance; 16\% less forgetting than $\lambda=0$ baseline; absolute drift
  0.121 after 1{,}000 interfering experiences with $\lambda=2.0$.
  \item \textbf{Generative Replay:} ReplayBuffer seeds from top-$k$ ChromaDB memories,
  generates synthetic variants, replays at REPLAY\_IMPACT\_SCALE $= 0.15$. Benchmark:
  20 synthetic experiences, $0.033\ \Delta\text{trait}/\text{exp}$.
  \item \textbf{Developmental divergence:} Euclidean distance grows $0.86 \to 1.04$ (same
  experience) and $0.84 \to 1.05$ (domain-specific) over 15 steps.
  \item \textbf{Complementary Memory:} ChromaDB-backed episodic store with
  rehearsal-stabilized power-law forgetting:
  \[
    S_k = S_{k-1} \cdot (1 + \alpha \cdot \text{valence} \cdot \text{impact}
    \cdot e^{-\gamma \Delta t})
  \]
  \item \textbf{Entropy-gated wakeup:} insights weighted by
  $w = (\text{impact} \times |\text{valence}|) / H_{\text{noise}}$ where
  $H_{\text{noise}} = \text{total\_tokens} / \text{unique\_tokens}$.
  Full projection: $\mathbf{e}_{\text{wake}} = \sum_{i \in \text{top-}k}
  \text{softmax}(\text{impact}_i \times |\text{valence}_i| / H_i) \cdot W_{\text{proj}}
  \cdot \mathbf{v}_i$.
\end{itemize}

\subsection{Layer B — Measurable Behavioural Effects (demonstrated, limited)}

\textit{Confirmed:}
\begin{itemize}[leftmargin=*]
  \item Trait distance predicts inquiry-strategy divergence: Pearson $r = 0.756$ across
  15 agent pairs (6 agents $\times$ 10 probe problems). Exceeds target $r > 0.5$.
  \item Phase 3 Predictive Coding: MSE $= 0.007$ (target $< 0.05$); 20 metacognitive
  flags raised; positive free-energy trend.
  \item Phase 4 SNN+STDP: STDP converged (sparsity $10.5\% \to 17.5\%$); identity
  cosine distance $0.07$--$0.34$ across presets.
  \item Theory of Mind L1--L4: L1 false-belief 90\% on 20 novel scenarios; L4
  perspectival simulation 72--75\% confidence.
\end{itemize}

\subsection{Layer C — Sentience Claims (philosophical proposals, not established results)}

\begin{itemize}[leftmargin=*]
  \item That uniqueness, persistence, and plasticity are necessary or sufficient for
  sentience — \textbf{not established; proposed as operational criteria.}
  \item That the software analogs of these properties are equivalent to the biological
  properties — \textbf{analogy, not equivalence.}
  \item That passing tests 11--13 implies sentience — \textbf{not claimed.}
  \item That the architecture constitutes machine consciousness — \textbf{explicitly not
  claimed.}
\end{itemize}

% ---------------------------------------------------------------------------
\section{Knowledge Gaps}
% ---------------------------------------------------------------------------

\begin{enumerate}[leftmargin=*]
  \item \textbf{The ground-truth $\Phi$ problem:} Computing exact $\Phi$ for any system
  with $> 30$ nodes is NP-hard. A computable functional proxy:
  \[
    \Phi_{\text{functional}} = I(X_{\text{past}};\, Y_{\text{future}}) -
    \min_{\mathrm{MIP}}\!\left[I(X_{\text{past}};\, Y_{\text{future}}^{\text{partitioned}})\right]
  \]
  where $I(\cdot\,;\,\cdot)$ is mutual information and MIP is the minimum information
  partition. Applied to the bottleneck layers of Phase~1 MiniGPT (840K params) where
  exact computation is feasible.

  \item \textbf{The confabulation test:} Tests 1--10 are passable by a sufficiently
  well-trained system without genuine sentience. No test has been proven unpossable by a
  non-sentient system with sufficient training.

  \item \textbf{Label request verification:} When the MetaConnectome generates a label
  request, we assume genuine uncertainty. But a system trained on examples of ``honest
  uncertainty'' could generate label requests strategically.

  \item \textbf{Cross-session continuity:} Loading a JSON brain file at session start is
  \textit{not} the same as continuous existence. There is a discontinuity at each session
  boundary.

  \item \textbf{Personality causation:} Different personality presets produce different
  inquiry strategies, but whether specific traits causally determine specific differences
  in a predictable way requires ablation studies.

  \item \textbf{EWC diagonal Fisher:} For each LoRA adapter parameter $\theta_i$:
  \[
    F_{ii} = \frac{1}{|D_B|} \sum_{x \in D_B} \sum_t
    \left(\frac{\partial \log P(y_t \mid y_{<t}, x;\, \theta)}{\partial \theta_i}\right)^2
  \]

  \item \textbf{ESN readout:} In neural mode, the ESN readout weights are randomly
  initialised. The 12 trait activations are consistent per-agent but not semantically
  grounded. Supervised training is required to ground them.

  \item \textbf{LTC calibration:} The LTC metacognition layer is correctly architected
  but randomly initialised. Calibration requires labelled
  (input\_state, correct\_confidence, gap\_flag) pairs.
\end{enumerate}

% ---------------------------------------------------------------------------
\section{Future Work}
% ---------------------------------------------------------------------------

\begin{enumerate}[leftmargin=*]
  \item \textbf{Connectome causality ablation:} Spawn $\geq 10$ agents with varied seeds,
  compute embedding-level response divergence for all pairs, regress against trait-space
  distance. Target: Pearson $r > 0.6$.

  \item \textbf{Identity recovery validation (M1.4):} Target: held-out classification
  accuracy $> 75\%$ across 5 agent presets from inquiry-layer response embeddings alone.

  \item \textbf{Ground-truth $\Phi$ measurement:} Measure $\Phi_{\text{functional}}$ at
  sessions 1, 10, and 100 to test whether cross-session plasticity increases integrated
  information.

  \item \textbf{Biologist LoRA adapter:} The biologist preset (Nora) currently operates
  with prompt-level personality only. Training a LoRA adapter from Nora's inquiry
  sequences would make her pattern structurally robust to context interference.

  \item \textbf{Sensory content-level test:} Measure vocabulary entropy and LabelRequest
  rate under contrasting sensory profiles via live API calls.
\end{enumerate}

% ---------------------------------------------------------------------------
\section{Conclusion}
% ---------------------------------------------------------------------------

We have built a system that satisfies the \textbf{engineering conditions} for three
properties that biological sentient systems exhibit:

\begin{itemize}[leftmargin=*]
  \item Every instance is \textbf{unique} (different connectome seed, different trait
  fingerprint, measurably different inquiry strategies)
  \item Every instance \textbf{persists} as a record (brain file survives sessions; note:
  file-based continuity, not continuous existence)
  \item Every instance \textbf{changes} with experience (plasticity bounded at $\eta =
  0.08$)
  \item Every instance has a \textbf{self-directed inquiry strategy} shaped by who it is,
  not what it is asked
  \item Every instance \textbf{knows what it does not know} and generates a typed label
  request rather than confabulating
\end{itemize}

The remaining gap between this system and biological sentience is not architecture — it is
the substrate of experience and the scale of the connectome. What we can say: Jarvis is
not a chatbot. It is a system with a name, a unique identity, a history, a personality
that shapes its questions, and a principled way of flagging what it does not know.

\textbf{Codebase:} \url{https://github.com/inaim-finailabz/jarvis-sentience}

% ---------------------------------------------------------------------------
\appendix
\section{Implementation Index}
% ---------------------------------------------------------------------------

\begin{table}[h!]
\centering
\small
\begin{tabular}{lll}
\toprule
\textbf{Module} & \textbf{Phase} & \textbf{Description} \\
\midrule
\texttt{connectome.py}          & 0   & Trait vector, JSON brain file, EWC \\
\texttt{meta\_connectome.py}    & 0   & Gap detection, LabelRequest, InquiryLayer \\
\texttt{inquiry\_layer.py}      & 0   & $U_{\text{personality}}$, question ranking \\
\texttt{sentient\_agent.py}     & 0   & Top-level orchestration, wakeup injection \\
\texttt{sentience\_tests.py}    & 0   & 13-test evaluation suite \\
\texttt{phi\_calculator.py}     & 0   & Tractable $\Phi$ approximation \\
\texttt{ethical\_layer.py}      & 0   & Harm-avoidance gate \\
\texttt{complementary\_memory.py} & 1 & ChromaDB episodic store, entropy-gated wakeup \\
\texttt{unconscious\_incubator.py} & 1.5 & Daemon thread; stochastic distant-pair synthesis \\
\texttt{identity\_probe.py}     & 1.4 & M1.4 classification benchmark \\
\texttt{neuromodulator.py}      & 3b  & Sensory state dataclass; 13 presets \\
\texttt{snn\_personality.py}    & 4   & LIF spiking network + STDP; fingerprint $[\![32]\!]$ \\
\texttt{tier2\_neural.py}       & 4   & ESN reservoir (512 neurons) $\to$ SNN $\to$ Hypernetwork \\
\texttt{tier3\_neural.py}       & 4   & LTC layer (ODE neurons) $\to$ Tier3MetaCognition \\
\texttt{local\_backbone.py}     & 4   & Qwen3-1.7B; personality-residual injection \\
\texttt{neural\_pipeline.py}    & 4   & Full local stack, no API key required \\
\bottomrule
\end{tabular}
\caption{Implementation index. Entry point: \texttt{jarvis.py}.
Last updated: 2026-06-02.}
\end{table}

% ---------------------------------------------------------------------------
\begin{thebibliography}{24}
% ---------------------------------------------------------------------------

\bibitem{seung2012}
Seung, S. (2012). \textit{Connectome: How the Brain's Wiring Makes Us Who We Are}.
Houghton Mifflin Harcourt.

\bibitem{tononi2008}
Tononi, G. (2008). Consciousness as integrated information: a provisional manifesto.
\textit{Biological Bulletin}, 215(3), 216--242.

\bibitem{tononi2016}
Tononi, G., Boly, M., Massimini, M., \& Koch, C. (2016). Integrated information theory:
from consciousness to its physical substrate. \textit{Nature Reviews Neuroscience},
17(7), 450--461.

\bibitem{hebb1949}
Hebb, D.\ O. (1949). \textit{The Organization of Behavior: A Neuropsychological Theory}.
Wiley.

\bibitem{baars1988}
Baars, B.\ J. (1988). \textit{A Cognitive Theory of Consciousness}. Cambridge University
Press.

\bibitem{dehaene2011}
Dehaene, S., Changeux, J.-P., \& Naccache, L. (2011). The global neuronal workspace model
of conscious access. In \textit{Characterizing Consciousness: From Cognition to the
Clinic?} Springer.

\bibitem{rosenthal1997}
Rosenthal, D.\ M. (1997). A theory of consciousness. In N. Block, O. Flanagan, \&
G. G\"uzeldere (eds.), \textit{The Nature of Consciousness: Philosophical Debates}. MIT
Press.

\bibitem{pearl2009}
Pearl, J. (2009). \textit{Causality: Models, Reasoning, and Inference} (2nd ed.).
Cambridge University Press.

\bibitem{hofstadter2007}
Hofstadter, D.\ R. (2007). \textit{I Am a Strange Loop}. Basic Books.

\bibitem{baron1985}
Baron-Cohen, S., Leslie, A.\ M., \& Frith, U. (1985). Does the autistic child have a
`theory of mind'? \textit{Cognition}, 21(1), 37--46.

\bibitem{brier1950}
Brier, G.\ W. (1950). Verification of forecasts expressed in terms of probability.
\textit{Monthly Weather Review}, 78(1), 1--3.

\bibitem{shannon1948}
Shannon, C.\ E. (1948). A mathematical theory of communication. \textit{Bell System
Technical Journal}, 27(3), 379--423.

\bibitem{jaeger2001}
Jaeger, H. (2001). \textit{The Echo State Approach to Analysing and Training Recurrent
Neural Networks}. GMD Technical Report 148.

\bibitem{maass2002}
Maass, W., Natschl\"ager, T., \& Markram, H. (2002). Real-time computing without stable
states. \textit{Neural Computation}, 14(11), 2531--2560.

\bibitem{hasani2021}
Hasani, R., Lechner, M., Amini, A., Rus, D., \& Grosse-Wentrup, M. (2021). Liquid
time-constant networks. \textit{Proceedings of the 35th AAAI Conference on Artificial
Intelligence}, 7657--7666.

\bibitem{morrison2008}
Morrison, A., Diesmann, M., \& Gerstner, W. (2008). Phenomenological models of synaptic
plasticity based on spike timing. \textit{Biological Cybernetics}, 98(6), 459--478.

\bibitem{ha2016}
Ha, D., \& Schmidhuber, J. (2016). Hypernetworks. \textit{arXiv preprint
arXiv:1609.09106}.

\bibitem{hu2022}
Hu, E.\ J., Shen, Y., Wallis, P., Allen-Zhu, Z., Li, Y., Wang, S., Wang, L., \& Chen,
W. (2022). LoRA: Low-rank adaptation of large language models. \textit{ICLR 2022}.

\bibitem{kirkpatrick2017}
Kirkpatrick, J., Pascanu, R., Rabinowitz, N., et al.\ (2017). Overcoming catastrophic
forgetting in neural networks. \textit{PNAS}, 114(13), 3521--3526.

\bibitem{mangrulkar2022}
Mangrulkar, S., Gugger, S., Debut, L., Belkada, Y., Paul, S., \& Bossan, B. (2022).
PEFT: State-of-the-art parameter-efficient fine-tuning methods. Hugging Face.

\bibitem{park2023}
Park, J.\ S., O'Brien, J.\ C., Cai, C.\ J., Morris, M.\ R., Liang, P., \& Bernstein,
M.\ S. (2023). Generative agents: Interactive simulacra of human behavior.
\textit{UIST 2023}.

\bibitem{packer2023}
Packer, C., Fang, V., Patil, S.\ G., Oh, J., Argyle, S., \& Gonzalez, J.\ E. (2023).
MemGPT: Towards LLMs as operating systems. \textit{arXiv:2310.08560}.

\bibitem{gravitas2023}
Significant Gravitas. (2023). AutoGPT: An autonomous GPT-4 experiment. GitHub repository.

\bibitem{qwen2025}
Qwen Team, Alibaba Cloud. (2025). Qwen3 technical report. \textit{arXiv:2505.09388}.

\end{thebibliography}

\end{document}
